\section{Memory Lanes and Theoretical Properties}
\label{sec:memory-lanes}
\label{sec:theory}

This section specifies the mathematical object implemented by the current
\ALUCLU reference backend.  The specification is intentionally narrower than
an abstract family of possible models: every state variable, update order, and
stabilizer below has a direct counterpart in
\texttt{src/aluclu}.  The implementation is token recurrent and uses explicit
caller-owned state.  It is therefore useful to distinguish three kinds of
statement throughout the section:
\begin{enumerate}
  \item algebraic identities of the real-arithmetic recurrence;
  \item conditional identities whose hypotheses state when numerical clamps
        are inactive; and
  \item empirical invariants exercised by the finite-precision test suite.
\end{enumerate}
In particular, ``exact'' denotes the absence of rank truncation for a retained
projected key--value record or live factor.  It does not denote unbounded,
lossless recall.

\subsection{Notation and block-level execution}
\label{subsec:notation-execution}

Let \(B\) be the batch size, \(T\) the number of tokens processed by one
call, \(d\) the model width, and \(L\) the number of residual blocks.
Table~\ref{tab:mathematical-symbols} maps the remaining symbols to the
configuration fields used by the implementation.

\begin{longtable}{@{}p{0.13\textwidth}p{0.36\textwidth}
  p{0.43\textwidth}@{}}
\caption{Symbols and their implementation counterparts.}
\label{tab:mathematical-symbols}\\
\toprule
Symbol & Code field & Meaning \\
\midrule
\endfirsthead
\toprule
Symbol & Code field & Meaning \\
\midrule
\endhead
\(H,d_k,d_v\)
  & \texttt{n\_heads}, \texttt{head\_key\_dim},
    \texttt{head\_value\_dim}
  & Number of Gated Delta heads and per-head key/value widths. \\
\(W\) & \texttt{local\_window}
  & Capacity of each local-attention rolling key--value window. \\
\(H_x,d_x,C_x\)
  & \texttt{ExactCacheConfig.n\_heads}, \(d/H_x\),
    \texttt{capacity}
  & Exact-cache head count, head width, and number of retained records. \\
\(S,K,C,A,r\)
  & \texttt{segment\_size}, \texttt{live\_segments},
    \texttt{archive\_slots}, \texttt{archive\_segments\_per\_slot},
    \texttt{archive\_rank}
  & Episodic segment length, exact-live count, archive-capsule count,
    segments per capsule, and archive rank. \\
\(d_r,p,m,J\)
  & \texttt{router\_dimension}, \texttt{router\_degree},
    \texttt{sketch\_dim}, \texttt{n\_sketches}
  & Route width, TensorSketch degree, sketch width, and number of sketches. \\
\(\kappa_g,\kappa_e\)
  & \path{GatedDeltaConfig.conv_kernel},
    \path{EpisodicMemoryConfig.conv_kernel}
  & Causal depthwise-stem kernel sizes. \\
\(f_{\min},f_{\max},\varepsilon\)
  & \texttt{feature\_floor}, \texttt{feature\_ceiling}, \texttt{eps}
  & Episodic feature bounds and numerical lower clamp. \\
\(b_x\) & \texttt{x.element\_size()}
  & Bytes per model-state scalar. \\
\bottomrule
\end{longtable}

For block \(\ell\), let \(x_{\ell,t}\in\R^d\) denote its input at token
\(t\), and define
\begin{equation}
  h_{\ell,t}=\RMSNorm(x_{\ell,t}).
  \label{eq:block-normalized-input}
\end{equation}
Every block contains a Gated Delta lane.  The local, episodic, and
residual-surprise lanes are present when, respectively,
\begin{align}
  \ell &\equiv 0
  \pmod{\texttt{local\_attention\_every}}, \\
  \ell &\equiv 0
  \pmod{\texttt{episodic\_every}}, \\
  \ell &\equiv 0
  \pmod{\texttt{exact\_cache\_every}},
\end{align}
using one-based layer indices, with the additional requirement that the
episodic or exact-cache configuration exists.  Consequently,
\begin{align}
  N_{\mathrm{local}}
    &=\left\lfloor
        \frac{L}{\texttt{local\_attention\_every}}
      \right\rfloor, \\
  N_{\mathrm{epi}}
    &=\mathbf{1}_{\{\text{episodic configured}\}}
      \left\lfloor
        \frac{L}{\texttt{episodic\_every}}
      \right\rfloor, \\
  N_{\mathrm{exact}}
    &=\mathbf{1}_{\{\text{exact cache configured}\}}
      \left\lfloor
        \frac{L}{\texttt{exact\_cache\_every}}
      \right\rfloor .
  \label{eq:scheduled-lane-counts}
\end{align}

All active lanes read the same \(h_{\ell,t}\); their outputs are not chained.
The Gated Delta lane additionally emits one scalar write-surprise value per
batch item and token.  The exact-cache lane consumes this scalar for its hard
insertion decision, but it computes its own query, key, and value projections.
Each memory lane performs its state update before reading at the same token.
Thus, the implemented convention is inclusive causal, \(j\leq t\).

\subsection{Causal depthwise streaming stems}
\label{subsec:streaming-stem}

The Gated Delta and episodic lanes each own an independent
\texttt{StreamingDepthwiseConv1d}.  For kernel width \(\kappa\), input chunk
\(x_{1:T}\), and incoming history \(c\in\R^{B\times(\kappa-1)\times d}\),
the operation is
\begin{align}
  \bar x &= [c;x_{1:T}], \\
  \widetilde x_{1:T}
    &=\operatorname{DepthwiseConv1d}_{\kappa}(\bar x), \\
  c'&=\bar x_{-(\kappa-1):}.
  \label{eq:streaming-convolution}
\end{align}
The history therefore stores exactly \(B(\kappa-1)d\) model-dtype elements.
The convolution is initialized as an identity map: all weights are zero
except the most recent causal tap, which is one, and the bias is zero.

\begin{proposition}[Chunk compositionality of the causal stem]
\label{prop:stem-compositionality}
In exact arithmetic, processing a sequence by
\eqref{eq:streaming-convolution} in any ordered partition produces the same
output sequence and terminal history as processing it in one call.
\end{proposition}
\noindent\emph{Proof sketch.}
At every boundary the state \(c'\) is exactly the suffix of length
\(\kappa-1\) required by the next causal convolution.  Each output at time
\(t\) is therefore evaluated on the same ordered \(\kappa\)-token context
under every partition.  Induction over chunk boundaries proves equality.
\hfill\(\square\)

\subsection{Gated Delta Rule-2 lane}
\label{subsec:gated-delta}

Suppress the block, batch, and head indices temporarily.  The lane applies its
own streaming stem to \(h_t\), producing \(\widetilde h_t\), and then forms
\begin{align}
  q_t
    &=\frac{W_q\widetilde h_t}
            {\max(\norm{W_q\widetilde h_t}_2,\varepsilon)}
      \in\R^{d_k}, \\
  k_t
    &=\frac{W_k\widetilde h_t}
            {\max(\norm{W_k\widetilde h_t}_2,\varepsilon)}
      \in\R^{d_k}, \\
  v_t&=W_v\widetilde h_t\in\R^{d_v}, \\
  e_t&=\sigmoid(W_e\widetilde h_t+b_e)\in(0,1)^{d_k}, \\
  w_t&=\sigmoid(W_w\widetilde h_t+b_w)\in(0,1)^{d_v}.
  \label{eq:gdr-projections}
\end{align}
Here \(e_t\) is the erase gate and \(w_t\) is the write gate.  Let
\(a=\exp(\texttt{decay\_log\_scale})>0\).  The token-dependent decay is
\begin{align}
  \log\gamma_t
    &=-a\odot
      \softplus\!\left(
        W_\gamma\widetilde h_t+\texttt{decay\_offset}
      \right), \\
  \gamma_t
    &=\max\!\left\{
        \exp(\log\gamma_t),\texttt{decay\_min}
      \right\},
  \label{eq:gdr-decay}
\end{align}
where the maximum is elementwise.  In real arithmetic,
\(\texttt{decay\_min}\leq\gamma_{t,j}<1\).

For each head, the persistent fast-weight matrix is
\(F_{t-1}\in\R^{d_k\times d_v}\).  The implemented recurrence is
\begin{align}
  \bar F_t
    &=\diag(\gamma_t)F_{t-1}, \\
  a_t&=e_t\odot k_t, \\
  \widehat v_t
    &=\bar F_t^\top a_t, \\
  v_t^{\mathrm{target}}
    &=w_t\odot v_t, \\
  \Delta v_t
    &=v_t^{\mathrm{target}}-\widehat v_t, \\
  F_t
    &=\bar F_t+k_t(\Delta v_t)^\top, \\
  o_t&=F_t^\top q_t.
  \label{eq:gdr-recurrence}
\end{align}
The outer-product update and read use the newly written \(F_t\), rather than
\(F_{t-1}\).  With head indices restored, the write-surprise scalar is
\begin{equation}
  \xi_t
  =
  \frac{
    \left(
      \sum_{h=1}^{H}\sum_{j=1}^{d_v}
      (\Delta v_{t,h,j})^2
    \right)^{1/2}
  }{\sqrt{Hd_v}}
  =
  \sqrt{
    \operatorname{mean}_{h,j}
    [(\Delta v_{t,h,j})^2]
  }.
  \label{eq:write-surprise}
\end{equation}
Thus \(\xi_t\) is an RMS write residual, not an unnormalized Euclidean norm.
The per-head output is gated and projected:
\begin{equation}
  g_t
  =
  W_{\mathrm{out}}
  \operatorname{concat}_{h=1}^{H}
  \left[
    o_{t,h}\odot
    \SiLU(W_o\widetilde h_t)_h
  \right].
  \label{eq:gdr-output}
\end{equation}
The implementation evaluates the fast-weight state, decay path, and recurrence
in fp32 even under mixed precision, then casts the stacked recurrent output
back to the model dtype before the output projection.

\begin{proposition}[Fixed recurrent state and inclusive causality]
\label{prop:gdr-state-causality}
For fixed \(H,d_k,d_v,\kappa_g,d\), the Gated Delta lane carries
\begin{equation}
  (\kappa_g-1)d
  \quad\text{model-dtype elements and}\quad
  Hd_kd_v
  \quad\text{fp32 elements}
  \label{eq:gdr-state-elements}
\end{equation}
per batch item, independently of the number of previously processed tokens.
Its output at \(t\) is a function only of inputs through \(t\), and it includes
the update induced by token \(t\).
\end{proposition}
\noindent\emph{Proof sketch.}
The only persistent tensors are the bounded convolution history and one
\(d_k\times d_v\) fast-weight matrix per head.  Every quantity on the
right-hand side of \eqref{eq:gdr-recurrence} depends on \(F_{t-1}\) and the
current stem output.  The read uses \(F_t\), establishing write-then-read
causality. \hfill\(\square\)

\subsection{Exact bounded local attention}
\label{subsec:local-attention}

Let \(H_\ell\) be the number of local-attention heads and
\(d_h=d/H_\ell\).  At token \(t\), projected \(k_t,v_t\) are appended before
the read and the cache is truncated to its last \(W\) entries:
\begin{equation}
  K_t=[K_{t-1};k_t]_{-W:},
  \qquad
  V_t=[V_{t-1};v_t]_{-W:}.
  \label{eq:local-window-update}
\end{equation}
For a retained position \(j\) and head \(h\),
\begin{equation}
  s_{t,j,h}
  =
  \frac{q_{t,h}^{\top}k_{j,h}}{\sqrt{d_h}}
  -
  \softplus(\texttt{slope\_raw}_h)(t-j),
  \label{eq:local-score}
\end{equation}
where the implementation constructs \(t-j\) as the reverse distance within
the rolling window.  It then computes
\begin{align}
  \alpha_{t,:,h}
    &=\softmax(s_{t,:,h}), \\
  \ell_{t,h}
    &=\sum_{j\in\mathcal W_t}\alpha_{t,j,h}v_{j,h},
  \qquad
  \mathcal W_t=\{\max(1,t-W+1),\ldots,t\}.
  \label{eq:local-read}
\end{align}
The softmax is explicitly evaluated in fp32 and cast back to the model dtype.
Attention-weight dropout is applied only in training mode.

\begin{proposition}[Local exactness and bound]
\label{prop:local-exactness}
The local lane computes ordinary causal softmax attention, with the learned
linear recency bias in \eqref{eq:local-score}, over exactly the retained
window \(\mathcal W_t\).  At saturation its persistent key--value payload is
\begin{equation}
  F_{\mathrm{local}}=2dW
  \label{eq:local-state-elements}
\end{equation}
model-dtype elements per batch item, and its attention read costs
\(\bigO(dW)\) per token, excluding dense projections.
\end{proposition}
\noindent\emph{Proof sketch.}
Equation~\eqref{eq:local-window-update} retains no index earlier than
\(t-W+1\), and it inserts the current key and value before
\eqref{eq:local-read}.  The softmax is not kernelized or low-rank
approximated.  Two tensors of shape \(H_\ell\times W\times d_h\) yield
\(2dW\) elements. \hfill\(\square\)

\subsection{Residual-surprise exact key--value cache}
\label{subsec:exact-cache}

The exact cache has independent projections
\begin{equation}
  q_t^x=W_q^x h_t,\qquad
  k_t^x=W_k^x h_t,\qquad
  v_t^x=W_v^x h_t,
  \label{eq:exact-cache-projections}
\end{equation}
split into \(H_x\) heads of width \(d_x=d/H_x\).  It receives the scalar
\(\xi_t\) from \eqref{eq:write-surprise}; no separate learned priority
projection is present.

\subsubsection{Hard insertion rule}

For each batch item, the state contains \(C_x\) key slots, value slots,
priorities, and valid bits.  Let
\begin{equation}
  \mathcal E_t
  =
  \{j:\neg\texttt{valid}_{t-1,j}\}.
\end{equation}
If \(\mathcal E_t\neq\varnothing\), the implementation selects the first
invalid slot.  Otherwise it selects
\begin{equation}
  j_{\min}
  =
  \arg\min_{1\leq j\leq C_x}
  \texttt{priority}_{t-1,j}.
\end{equation}
The comparison priority is detached from autograd and converted to the
priority dtype:
\begin{equation}
  \widehat\xi_t=\stopgrad(\xi_t).
\end{equation}
The write predicate is
\begin{equation}
  \operatorname{write}_t
  =
  \mathbf{1}\{\mathcal E_t\neq\varnothing\}
  \lor
  \mathbf{1}\{
    \widehat\xi_t>
    \texttt{priority}_{t-1,j_{\min}}
  \}.
  \label{eq:exact-cache-write}
\end{equation}
The inequality is strict: an equal-priority arrival does not replace the
older retained record.  All heads write the current projected key and value to
the same selected slot.

\begin{proposition}[Stable top-\(C_x\) priority invariant]
\label{prop:top-cx}
Assume finite priorities and process one batch item in temporal order.  After
\(t\) updates, the cache contains \(\min(C_x,t)\) records whose priorities are
the largest priorities seen so far, with earlier records retained at the
cutoff when priorities tie.
\end{proposition}
\noindent\emph{Proof sketch.}
Before saturation every arrival occupies one empty position.  After
saturation, an arrival below or equal to the current minimum cannot belong to
a strictly better top-\(C_x\) set and is rejected.  An arrival above the
minimum replaces exactly one minimum.  Induction on \(t\) gives the invariant;
strict comparison yields stable earlier-tie retention. \hfill\(\square\)

\subsubsection{Write-then-read exact softmax}

Insertion precedes the current read.  For valid slot \(j\),
\begin{align}
  s_{t,h,j}^x
    &=
    \frac{(q_{t,h}^x)^\top k_{j,h}^x}{\sqrt{d_x}}, \\
  \alpha_{t,h,:}^x
    &=\softmax(s_{t,h,:}^x), \\
  o_{t,h}^x
    &=\sum_{j:\mathrm{valid}_j}
      \alpha_{t,h,j}^x v_{j,h}^x.
  \label{eq:exact-cache-read}
\end{align}
Invalid positions are masked with the minimum finite value of the score dtype.
The first token is always inserted because capacity is positive and an empty
slot exists, so every read has at least one valid entry.  Half and bfloat16
scores use fp32 softmax; float32 and float64 scores retain their score dtype.
The current record participates in its read if accepted.  If rejected, its
query still reads the retained past, but its key and value are absent.

The selection path is deliberately non-differentiable:
\(\widehat\xi_t\) is detached, and integer indexing plus one-hot selection
implements the replacement.  Accepted key/value projections receive gradient
through \eqref{eq:exact-cache-read}; the cache loss does not backpropagate
through the discrete ranking decision into \(\xi_t\).

\begin{remark}[Meaning of exactness]
The cache stores the projected \(k_t^x,v_t^x\) of accepted records without
rank truncation or averaging.  It does not store all tokens, raw hidden states,
token identifiers, or any record rejected or later evicted by
\eqref{eq:exact-cache-write}.
\end{remark}

Its logical state size per batch item is
\begin{equation}
  F_{\mathrm{exact}}
  =
  2C_xd+2C_x.
  \label{eq:exact-cache-state-elements}
\end{equation}
For ordinary fp16, bfloat16, or fp32 model state, priority is fp32 and valid is
Boolean, so the byte payload is
\begin{equation}
  \operatorname{bytes}_{\mathrm{exact},B=1}
  =
  2b_xC_xd+5C_x.
  \label{eq:exact-cache-state-bytes}
\end{equation}
In the float64 validation path priority is also float64, replacing \(5C_x\)
by \(9C_x\).  Full-shape \texttt{torch.where} updates and the softmax read
both cost \(\bigO(C_xd)\) per token, excluding dense projections.

\subsection{Bounded episodic associative memory}
\label{subsec:episodic-memory}

The episodic lane applies its own causal stem to \(h_t\), yielding
\(\widetilde h_t^e\), and uses a bounded positive feature map
\begin{equation}
  \phi(a)
  =
  f_{\min}+(f_{\max}-f_{\min})\sigmoid(a),
  \qquad
  0<f_{\min}<f_{\max}.
  \label{eq:positive-feature-map}
\end{equation}
The implementation computes
\begin{align}
  q_t^+&=\phi(W_q^e\widetilde h_t^e), \\
  k_t^+&=\phi(W_k^e\widetilde h_t^e), \\
  v_t^e&=W_v^e\widetilde h_t^e.
  \label{eq:episodic-projections}
\end{align}
Componentwise,
\begin{equation}
  f_{\min}\leq \phi(a)_j\leq f_{\max},
  \label{eq:feature-bounds}
\end{equation}
where equality can occur through floating-point sigmoid saturation.

\subsubsection{Separate numerator and denominator summaries}

For a retained token set \(\mathcal I_i\), slot \(i\) represents
\begin{align}
  M_i
    &=\sum_{t\in\mathcal I_i}
      v_t^e(k_t^+)^\top, \\
  z_i
    &=\sum_{t\in\mathcal I_i}k_t^+.
  \label{eq:slot-moments}
\end{align}
The numerator is stored in factors
\begin{equation}
  M_i
  =
  L_i^\top R_i,
  \qquad
  L_i,R_i\in\R^{f_i\times d},
  \label{eq:slot-factorization}
\end{equation}
corresponding to the code fields \texttt{left} and \texttt{right}.  For an
exact current or live segment, one factor row corresponds to one token.  The
slot read is
\begin{align}
  n_i(q_t^+)
    &=M_iq_t^+
      =L_i^\top(R_iq_t^+), \\
  d_i(q_t^+)
    &=z_i^\top q_t^+.
  \label{eq:slot-read}
\end{align}
The denominator is never reconstructed from \(R_i\).  After SVD truncation,
\(R_i\) is a signed low-rank factor and is not a collection of positive keys.
The independent \(z_i\) is therefore the sole valid denominator summary.

\begin{proposition}[Positive denominator and compression separation]
\label{prop:positive-denominator}
Every nonempty retained slot satisfies
\begin{equation}
  d_i(q_t^+)>0
\end{equation}
in real arithmetic.  Rank truncation of \(M_i\) changes only the numerator in
\eqref{eq:slot-read}; it does not change \(z_i\) or \(d_i\).
\end{proposition}
\noindent\emph{Proof sketch.}
Every component of every retained \(k_t^+\) and of \(q_t^+\) is positive by
\eqref{eq:positive-feature-map}.  Hence \(z_i\) is componentwise positive and
its inner product with \(q_t^+\) is positive.  The compression routines replace
only \texttt{left} and \texttt{right}; they copy or add \(z_i\) directly.
\hfill\(\square\)

\subsubsection{Route features and TensorSketch}

The route vector is
\begin{equation}
  u_t
  =
  \frac{W_r\widetilde h_t^e}
       {\max(\norm{W_r\widetilde h_t^e}_2,\varepsilon)}
  \in\R^{d_r}.
  \label{eq:route-vector}
\end{equation}
Thus \(\norm{u_t}_2=1\) when the pre-normalization norm is at least
\(\varepsilon\), and \(\norm{u_t}_2\leq 1\) otherwise.

For sketch \(j\in\{1,\ldots,J\}\), degree factor
\(a\in\{1,\ldots,p\}\), bucket \(b\in\{0,\ldots,m-1\}\), fixed hash
\(h_{j,a}\), and fixed sign map \(s_{j,a}\), define
\begin{equation}
  \operatorname{CS}_{j,a}(u)[b]
  =
  \sum_{\ell:h_{j,a}(\ell)=b}
  s_{j,a}(\ell)u_\ell .
  \label{eq:count-sketch}
\end{equation}
The implemented TensorSketch feature is
\begin{equation}
  \psi_j(u)
  =
  \operatorname{IFFT}\!\left(
    \prod_{a=1}^{p}
    \operatorname{FFT}(
      \operatorname{CS}_{j,a}(u)
    )
  \right)
  \in\R^m.
  \label{eq:tensor-sketch}
\end{equation}
Hashes and signs are generated once from \texttt{router\_seed} and registered
as persistent buffers.  A slot stores additive summaries
\begin{align}
  r_i&=\sum_{t\in\mathcal I_i}u_t, \\
  \Psi_{i,j}
    &=\sum_{t\in\mathcal I_i}\psi_j(u_t), \\
  c_i&=\abs{\mathcal I_i}.
  \label{eq:route-summaries}
\end{align}
The sketch is used only to score slots.  Finite \(m\) and \(J\) permit
collisions, and the implementation does not claim a collision-free identity
store.  Configuration validation requires odd \(J\), making the median below
unambiguous.

\subsubsection{Mass-weighted routing and its clamp condition}

For route query \(u_q\) and slot \(i\), the linear and polynomial scores are
\begin{align}
  \ell_i
    &=
    \frac{u_q^\top r_i}{\sqrt{\max(c_i,1)}}, \\
  \rho_i
    &=
    \frac{
      \median_{j=1,\ldots,J}
      [\psi_j(u_q)^\top\Psi_{i,j}]
    }{\sqrt{\max(c_i,1)}}.
  \label{eq:route-component-scores}
\end{align}
Let
\begin{equation}
  (\pi_1,\pi_p)=\softmax(\texttt{mix\_logits}),
  \qquad
  s_i=\pi_1\ell_i+\pi_p\rho_i.
  \label{eq:mixed-route-score}
\end{equation}
The signed, bounded temperature is
\begin{equation}
  \theta
  =
  \texttt{max\_temperature}
  \tanh(\texttt{temperature\_raw}).
  \label{eq:router-temperature}
\end{equation}
The implementation does not constrain \(\theta\) to be positive.

It is important to state the numerical stabilizers exactly.  With
\begin{align}
  D_0&=\sum_i d_i, \\
  D_\varepsilon&=\max(D_0,\varepsilon), \\
  \bar s
    &=\frac{\sum_i s_i d_i}{D_\varepsilon}, \\
  \eta_i
    &=\clip\!\left(
       \theta(s_i-\bar s),-L_w,L_w
      \right), \\
  \widetilde w_i&=\exp(\eta_i)>0, \\
  c_\varepsilon
    &=\max\!\left(
      \frac{\sum_i\widetilde w_i d_i}{D_\varepsilon},
      \varepsilon
    \right), \\
  w_i&=\frac{\widetilde w_i}{c_\varepsilon},
  \qquad
  L_w=\texttt{max\_log\_weight},
  \label{eq:implemented-router}
\end{align}
the routed numerator and returned denominator are
\begin{align}
  N_{\mathrm{route}}
    &=\sum_iw_i n_i, \\
  D_{\mathrm{returned}}
    &=D_0, \\
  y_t^e
    &=\frac{N_{\mathrm{route}}}
            {\max(D_{\mathrm{returned}},\varepsilon)}.
  \label{eq:episodic-output}
\end{align}
The code intentionally does not multiply the returned denominator by \(w_i\).

\begin{proposition}[Conditional denominator-mass conservation]
\label{prop:conditional-mass}
If both lower clamps in \eqref{eq:implemented-router} are inactive, namely
\begin{equation}
  D_0\geq\varepsilon
  \quad\text{and}\quad
  \frac{\sum_i\widetilde w_i d_i}{D_0}
  \geq\varepsilon,
  \label{eq:mass-conservation-conditions}
\end{equation}
then
\begin{equation}
  w_i>0,
  \qquad
  \sum_iw_id_i=D_0.
  \label{eq:mass-conservation}
\end{equation}
If additionally \(\theta=0\), then \(w_i=1\) for every slot.
\end{proposition}
\noindent\emph{Proof sketch.}
Under \eqref{eq:mass-conservation-conditions},
\[
  c_\varepsilon
  =
  \frac{\sum_i\widetilde w_i d_i}{D_0}.
\]
Substitution into \(w_i=\widetilde w_i/c_\varepsilon\) yields
\(\sum_iw_id_i=D_0\); positivity follows from the exponential and positive
normalizer.  When \(\theta=0\), every \(\eta_i=0\), hence
\(\widetilde w_i=1\) and \(c_\varepsilon=1\). \hfill\(\square\)

\begin{remark}[Active-clamp behavior]
\label{rem:active-clamp}
When either condition in \eqref{eq:mass-conservation-conditions} fails, the
implemented stabilized formula remains \eqref{eq:implemented-router}, but
\eqref{eq:mass-conservation} is not an unconditional identity.  For example,
if only the first clamp is active and the second is inactive, then
\(\sum_iw_id_i=D_\varepsilon\), not necessarily \(D_0\).  This distinction is
why all exact mass-conservation claims in this paper carry the inactive-clamp
hypothesis.  The unit test exercises positive masses well above
\(\varepsilon\) and verifies the equality to floating-point tolerance.
\end{remark}

\subsubsection{Current, live, and archive lifecycle}

Each token first creates a one-row slot and is appended to
\texttt{current}.  When the current factor count reaches \(S\), the completed
slot moves to the end of the exact \texttt{live} tuple and
\texttt{current} becomes empty.  If this produces more than \(K\) live
segments, the oldest live segment is evicted.  When \(C=0\), that segment is
forgotten.  Otherwise it enters a temporal archive ring:
\begin{enumerate}
  \item the first segment of a capsule is compressed to rank \(r\);
  \item further incoming segments are merged into the active capsule until it
        contains \(A\) segments;
  \item the cursor then advances to the next capsule; and
  \item once all \(C\) positions exist, advancing overwrites the next capsule
        and erases its former contents.
\end{enumerate}
The physical tuple order is not required to be chronological; the cursor
identifies the active ring position.

The method \texttt{maximum\_retained\_tokens()} returns the conservative
allocation-style horizon
\begin{equation}
  H_{\max}
  =
  S(1+K+CA)
  =
  \underbrace{S}_{\text{current budget}}
  +
  \underbrace{KS}_{\text{exact live}}
  +
  \underbrace{CAS}_{\text{compressed archive}}.
  \label{eq:retained-horizon}
\end{equation}
Because a current segment is sealed immediately at \(S\),
\begin{equation}
  0\leq\abs{\texttt{current}}<S,
  \qquad
  N_{\mathrm{retained}}\leq H_{\max}-1
  \label{eq:reachable-retained-horizon}
\end{equation}
for an actually reachable state.  After a full-ring overwrite, the number of
archived segments falls from \(CA\) to \((C-1)A+1\), then rises as the new
active capsule fills.  Hence the age horizon has a sawtooth variation of up to
\(A-1\) archived segments.

\begin{proposition}[Bounded episodic lifecycle]
\label{prop:bounded-episodic-lifecycle}
For every stream length,
\begin{equation}
  \abs{\texttt{live}}\leq K,\qquad
  \abs{\texttt{archive}}\leq C,\qquad
  \texttt{archive\_cursor\_segments}\leq A,
  \label{eq:episodic-container-bounds}
\end{equation}
and \eqref{eq:reachable-retained-horizon} holds.  Current and live numerator
factors are algebraically exact sums for their retained tokens; archive
numerators have rank at most \(r\); overwritten capsules have no remaining
representation.
\end{proposition}
\noindent\emph{Proof sketch.}
Sealing resets current before it can exceed \(S-1\).  Live overflow immediately
evicts its oldest element, maintaining \(K\).  Archive insertion either merges
into the current capsule, appends while fewer than \(C\) capsules exist, or
overwrites one existing capsule.  The cursor advances after \(A\) segments.
These transitions establish \eqref{eq:episodic-container-bounds} by induction;
counting their maximum represented tokens yields
\eqref{eq:reachable-retained-horizon}. \hfill\(\square\)

\subsection{Thin-QR, small-core SVD compression}
\label{subsec:small-core-svd}

The archive compresses a factor pair without materializing a dense
\(d\times d\) numerator.  Suppress the batch index and let
\begin{equation}
  L=\texttt{left}^{\top}\in\R^{d\times F},
  \qquad
  R=\texttt{right}^{\top}\in\R^{d\times F},
  \qquad
  M=LR^\top.
  \label{eq:compression-input}
\end{equation}
Reduced QR decompositions give
\begin{equation}
  L=Q_LR_L,\qquad
  R=Q_RR_R.
\end{equation}
The code performs an SVD only on
\begin{equation}
  C_{\mathrm{core}}
  =
  R_LR_R^\top
  =
  U_c\Sigma V_c^\top.
  \label{eq:small-core}
\end{equation}
If \(p_c=\min(d,F)\), then \(C_{\mathrm{core}}\in\R^{p_c\times p_c}\).
Let
\begin{equation}
  r_{\mathrm{eff}}
  =
  \min(r,\operatorname{len}(\Sigma)).
\end{equation}
The returned factors are
\begin{align}
  \widehat L_{\mathrm{rows}}
    &=
    \left(
      Q_LU_{c,:r_{\mathrm{eff}}}
      \Sigma_{:r_{\mathrm{eff}}}^{1/2}
    \right)^\top, \\
  \widehat R_{\mathrm{rows}}
    &=
    \left(
      Q_RV_{c,:r_{\mathrm{eff}}}
      \Sigma_{:r_{\mathrm{eff}}}^{1/2}
    \right)^\top.
  \label{eq:compressed-factors}
\end{align}
Zero rows pad each factor to exactly \(r\), keeping archive shapes static.
For half or bfloat16 input the QR and SVD work in fp32 before results are cast
back; float32 and float64 retain their input work dtype.

\begin{proposition}[Best rank-\(r\) numerator in exact arithmetic]
\label{prop:best-rank}
Ignoring finite-precision QR/SVD and the final dtype cast, the reconstructed
\begin{equation}
  \widehat M
  =
  \widehat L_{\mathrm{rows}}^\top
  \widehat R_{\mathrm{rows}}
\end{equation}
is a best rank-\(r_{\mathrm{eff}}\) approximation to \(M\) in Frobenius norm,
and
\begin{equation}
  \norm{M-\widehat M}_F
  =
  \delta
  =
  \left(
    \sum_{j>r_{\mathrm{eff}}}\sigma_j(C_{\mathrm{core}})^2
  \right)^{1/2}.
  \label{eq:discarded-tail}
\end{equation}
If \(r\geq\rank(M)\), the compression is algebraically exact.
\end{proposition}
\noindent\emph{Proof sketch.}
Reduced QR factors have orthonormal columns, so left and right multiplication
by \(Q_L,Q_R^\top\) preserves all nonzero singular values and Frobenius norms.
Truncating the SVD of \(C_{\mathrm{core}}\) therefore gives the
Eckart--Young best-rank approximation to \(M\).  Splitting
\(\Sigma^{1/2}\) symmetrically reconstructs that truncation.
\hfill\(\square\)

In a typical merge, an existing rank-\(r\) capsule is concatenated with an
exact \(S\)-row live segment, so \(F\leq r+S\) and
\begin{equation}
  p_c=\min(d,F)\leq\min(d,r+S).
\end{equation}
One merge costs approximately
\begin{equation}
  \bigO(dFp_c+p_c^3).
  \label{eq:compression-merge-cost}
\end{equation}
Since a live overflow occurs at most once every \(S\) tokens, its steady-state
amortized cost is
\begin{equation}
  \bigO\!\left(
    \frac{d(r+S)p_c+p_c^3}{S}
  \right)
  \quad\text{per token}.
  \label{eq:compression-amortized-cost}
\end{equation}
The core is genuinely small only when \(r+S\ll d\).  Configuration validation
enforces \(r\leq d\), but does not enforce \(r+S<d\); the core can therefore
grow to \(d\times d\).

\subsection{Compression error ledger}
\label{subsec:error-ledger}

For one archive capsule, let \(M_1,\ldots,M_u\) be exact incoming segment
numerators and
\begin{equation}
  M_u^\star=\sum_{j=1}^{u}M_j
\end{equation}
the ideal retained numerator.  At merge \(u\), the implementation truncates
\begin{equation}
  A_u=\widehat M_{u-1}+M_u,
  \qquad
  \widehat M_u=P_r(A_u),
\end{equation}
and reports
\begin{equation}
  \delta_u
  =
  \norm{A_u-P_r(A_u)}_F
  =
  \left(
    \sum_{j>r}\sigma_j(A_u)^2
  \right)^{1/2}.
  \label{eq:merge-tail}
\end{equation}
The slot field \texttt{error\_bound} is updated linearly:
\begin{equation}
  E_u
  =
  E_{u-1}+E_{\mathrm{incoming}}+\delta_u.
  \label{eq:error-ledger-update}
\end{equation}
For a freshly evicted exact live segment,
\(E_{\mathrm{incoming}}=0\).

\begin{proposition}[Retained-capsule error bound]
\label{prop:error-ledger}
In exact arithmetic, the ledger satisfies
\begin{equation}
  \norm{M_u^\star-\widehat M_u}_F\leq E_u.
  \label{eq:ledger-frobenius-bound}
\end{equation}
Consequently, for any query \(q\),
\begin{equation}
  \norm{(M_u^\star-\widehat M_u)q}_2
  \leq E_u\norm{q}_2.
  \label{eq:ledger-query-bound}
\end{equation}
\end{proposition}
\noindent\emph{Proof sketch.}
Write
\[
  M_u^\star-\widehat M_u
  =
  (M_{u-1}^\star-\widehat M_{u-1})
  +(M_u-\widehat M_{\mathrm{incoming}})
  +(A_u-P_r(A_u)).
\]
The triangle inequality and the inductive bounds on the first two terms yield
\eqref{eq:ledger-frobenius-bound} with the additive update
\eqref{eq:error-ledger-update}.  Submultiplicativity,
\(\norm{Eq}_2\leq\norm{E}_F\norm{q}_2\), yields
\eqref{eq:ledger-query-bound}. \hfill\(\square\)

Because \(z_i\), route sums, sketch sums, and counts are accumulated separately
from numerator compression, the ideal and compressed reads use the same
denominator and router weights when compared at the same retained segmentation.
For positive \(w_i\),
\begin{equation}
  \norm{y^\star-\widehat y}_2
  \leq
  \frac{
    \norm{q}_2
    \sum_{i\in\mathrm{archive}}w_iE_i
  }{D_0},
  \label{eq:routed-output-error-bound}
\end{equation}
provided the mathematical denominator is \(D_0>0\).  This bound covers only
currently retained capsules.  It excludes overwritten history, finite-precision
QR/SVD error, rounding from casting factors back to a lower dtype, and any task
accuracy or recall guarantee.

\subsection{Token-wise branch fusion and SwiGLU}
\label{subsec:fusion}

Let \(b_{t,j}\in\R^d\) denote active branch \(j\)'s output.  Each branch has an
independent RMS normalization and an unconstrained learned scalar:
\begin{equation}
  \widetilde b_{t,j}
  =
  \beta_j\RMSNorm_j(b_{t,j}),
  \qquad
  \beta_j=\texttt{branch\_scale}_j.
  \label{eq:branch-normalization}
\end{equation}
When more than one branch is active,
\begin{align}
  \alpha_t
    &=\softmax(W_fh_t+b_f), \\
  f_t
    &=\sum_j\alpha_{t,j}\widetilde b_{t,j},
  \qquad
  \alpha_{t,j}>0,\quad
  \sum_j\alpha_{t,j}=1.
  \label{eq:convex-fusion}
\end{align}
With a single branch, the fusion projection is absent and the normalized,
scaled branch is returned directly.  For a multi-branch block, \(W_f,b_f\)
are initialized to zero and every \(\beta_j\) to one, giving a uniform mixture
at initialization.

\begin{proposition}[Precise convexity statement]
\label{prop:fusion-convexity}
For each token, \(f_t\) lies in the convex hull of the scaled normalized
vectors \(\{\widetilde b_{t,j}\}_j\).  It is not generally constrained to the
convex hull of the raw branch outputs \(\{b_{t,j}\}_j\).
\end{proposition}
\noindent\emph{Proof sketch.}
Softmax gives positive coefficients summing to one, proving the first claim.
The second follows because each \(\beta_j\) is unconstrained and each branch
has a distinct nonlinear RMS normalization. \hfill\(\square\)

The first residual update is
\begin{equation}
  x'_t
  =
  x_t+\operatorname{Dropout}(f_t).
\end{equation}
Writing
\begin{equation}
  (g_t^{\mathrm{ffn}},v_t^{\mathrm{ffn}})
  =
  \operatorname{split}\!\left(
    W_{\mathrm{in}}\RMSNorm(x'_t)
  \right),
\end{equation}
the feed-forward output and second residual are
\begin{align}
  \operatorname{SwiGLU}(x'_t)
    &=
    W_{\mathrm{out}}^{\mathrm{ffn}}
    \left[
      \SiLU(g_t^{\mathrm{ffn}})
      \odot v_t^{\mathrm{ffn}}
    \right], \\
  x_{t,\mathrm{out}}
    &=
    x'_t+
    \operatorname{Dropout}(
      \operatorname{SwiGLU}(x'_t)
    ).
  \label{eq:block-output}
\end{align}

\subsection{Global causality and chunk/step equivalence}
\label{subsec:causality-chunking}

Let a lane state transition be written
\[
  (y_t,\state_t)=\mathcal T(x_t,\state_{t-1}).
\]
The concrete \(\state_t\) is the tuple of convolution histories, Gated Delta
fast weights, local key--value windows, exact-cache tensors and metadata, and
episodic current/live/archive state.

\begin{proposition}[Inclusive causal block semantics]
\label{prop:global-causality}
With fixed parameters and incoming state, the output at token \(t\) depends
only on the incoming state and tokens through \(t\).  No future token can
change an earlier output.  The current token can influence its own memory
read in every memory lane.
\end{proposition}
\noindent\emph{Proof sketch.}
The streaming stems are causal by
\eqref{eq:streaming-convolution}.  The Gated Delta lane updates \(F_t\) before
reading it; local attention appends \((k_t,v_t)\) before softmax; the exact
cache decides insertion before reading; and the episodic lane writes and
performs any seal, eviction, merge, or overwrite before aggregating its slots.
Fusion and SwiGLU are token local.  Composition over layers preserves
causality. \hfill\(\square\)

The language-model loss uses
\(\texttt{logits[:, :-1]}\) against \(\texttt{labels[:, 1:]}\), so the
inclusive memory read is compatible with next-token prediction.

\begin{proposition}[Chunk and step equivalence]
\label{prop:chunk-step}
Assume evaluation mode, fixed parameters, identical initial state, exact
arithmetic, and the same token order.  For any ordered partition of a sequence,
concatenating chunk outputs yields the same outputs and terminal state as one
full call or repeated one-token \texttt{step} calls.
\end{proposition}
\noindent\emph{Proof sketch.}
Each forward method applies the same token transition in temporal order.
Proposition~\ref{prop:stem-compositionality} handles both convolutional stems.
All other state updates are functions solely of the previous explicit state
and current projected token.  Associating the composition
\(\mathcal T_T\circ\cdots\circ\mathcal T_1\) according to chunk boundaries
does not change it. \hfill\(\square\)

\begin{remark}[Finite precision and training mode]
The implementation tests chunk/step equality in evaluation mode with numerical
tolerances, because vectorized linear projections and convolution kernels may
have shape-dependent floating-point reduction order.  In training mode,
different partitioning can also consume dropout random numbers differently.
Therefore Proposition~\ref{prop:chunk-step} is not stated as bitwise equality
for arbitrary training executions.
\end{remark}

\subsection{Persistent-state and compute bounds}
\label{subsec:state-compute}

Define the full conservative current/live slot size
\begin{equation}
  F_{\mathrm{live}}
  =
  2Sd+d+d_r+Jm+2,
  \label{eq:live-slot-size}
\end{equation}
and archive slot size
\begin{equation}
  F_{\mathrm{archive}}
  =
  2rd+d+d_r+Jm+2.
  \label{eq:archive-slot-size}
\end{equation}
The terms are the two numerator factors, exact \(z\), route sum, \(J\) sketch
sums, count, and error ledger.  The episodic state bound per batch item is
\begin{equation}
  F_{\mathrm{epi}}
  =
  (K+1)F_{\mathrm{live}}
  +CF_{\mathrm{archive}}
  +(\kappa_e-1)d.
  \label{eq:episodic-state-bound}
\end{equation}
This is conservative because a reachable current slot has fewer than \(S\)
rows.  The temporary fp32 QR/SVD workspace is not persistent state.

Combining all scheduled lanes gives the following logical tensor-byte bound
per batch item for ordinary non-float64 execution:
\begin{align}
  \operatorname{bytes}_{B=1}
  \leq{}&
  L\left[
    b_x(\kappa_g-1)d
    +4Hd_kd_v
  \right]
  +
  N_{\mathrm{local}}(2b_xdW)
  \nonumber\\
  &+
  N_{\mathrm{epi}}(b_xF_{\mathrm{epi}})
  +
  N_{\mathrm{exact}}
  \left(
    2b_xC_xd+5C_x
  \right).
  \label{eq:model-state-byte-bound}
\end{align}
Batch size \(B\) multiplies the tensor payload by \(B\).  This accounting
excludes parameters, optimizer state, autograd activations, QR/SVD workspace,
Python metadata, allocator fragmentation, and excess backing-storage
capacity.  In particular, bounded inference state does not imply bounded
training activation memory when the entire computation graph is retained.

Table~\ref{tab:lane-complexity} summarizes token costs.  Dense projection costs
are shown separately because the persistent-memory terms are the quantities
that distinguish the four lanes.

\begin{table}[t]
\centering
\small
\caption{Per-token time and persistent state for fixed configuration.}
\label{tab:lane-complexity}
\begin{tabularx}{\linewidth}{@{}l X X@{}}
\toprule
Lane & Per-token time beyond or including stated projections
     & Persistent elements per batch item \\
\midrule
Gated Delta
  & \(\bigO(Hd_kd_v+\kappa_gd)\) recurrence/stem, plus its dense
    query/key/value/gate projections
  & \(Hd_kd_v\) fp32 plus \((\kappa_g-1)d\) model-dtype \\
Local exact
  & \(\bigO(dW)\) attention plus \(\bigO(d^2)\) QKV/output projections
  & \(2dW\) model-dtype \\
Residual-surprise exact
  & \(\bigO(C_xd)\) selection, full-shape update, and read, plus
    \(\bigO(d^2)\) projections
  & \(2C_xd\) model-dtype plus \(C_x\) priority and \(C_x\) valid entries \\
Episodic
  & Factor read, routing, sketching, dense projections, and amortized
    compression as in \eqref{eq:episodic-token-cost}
  & At most \(F_{\mathrm{epi}}\) model-dtype elements \\
\bottomrule
\end{tabularx}
\end{table}

At most \(1+K+C\) episodic slots participate in a read.  Their factor-read
cost is
\begin{equation}
  \bigO\!\left(
    d[S(K+1)+Cr]
  \right).
\end{equation}
Route comparisons cost
\begin{equation}
  \bigO\!\left(
    (K+C+1)(d_r+Jm)
  \right),
\end{equation}
and the current FFT implementation constructs a token sketch in
\begin{equation}
  \bigO\!\left(
    Jp[d_r+m\log m]
  \right).
\end{equation}
Including amortized compression but suppressing dense projection constants,
\begin{equation}
  \begin{split}
  \operatorname{cost}_{\mathrm{epi/token}}
  =\bigO\bigg(
    &d[S(K+1)+Cr]
    +(K+C+1)(d_r+Jm)\\
    &+Jp[d_r+m\log m]
    +\frac{d(r+S)p_c+p_c^3}{S}
  \bigg).
  \end{split}
  \label{eq:episodic-token-cost}
\end{equation}
For a fixed configuration, recurrent-decode cost and persistent state are
independent of total context length.  A \(T\)-token scan nevertheless costs
\(\bigO(T)\) times the fixed per-token work.  The current Python backend uses
token loops, tuple concatenation, full-cache \texttt{torch.where}, and
occasional QR/SVD; these asymptotic state bounds do not imply fused-kernel
throughput.

\subsection{Finite-capacity boundary}
\label{subsec:finite-capacity}

Every persistent lane is explicitly bounded:
\begin{itemize}
  \item Gated Delta has \(Hd_kd_v\) fp32 fast-weight elements;
  \item local attention stores only the last \(W\) projected key--value pairs;
  \item the residual-surprise cache stores at most \(C_x\) selected projected
        key--value pairs; and
  \item episodic memory stores one partial segment, at most \(K\) exact live
        segments, and at most \(C\) rank-\(r\) archive capsules.
\end{itemize}
For every capsule,
\begin{equation}
  \rank(\widehat M_i)\leq r.
\end{equation}
Exact \(z_i\), route, sketch, and count summaries cannot reconstruct numerator
directions discarded by rank truncation.

\begin{proposition}[Finite-state impossibility of arbitrary unbounded recall]
\label{prop:finite-state-impossibility}
Fix model parameters and a persistent state representation with at most
\(Q\) physical bits.  Suppose query keys are supplied externally and each of
\(N\) independent associated values belongs to an alphabet
\(\mathcal V\), with \(\abs{\mathcal V}\geq2\).  A system that recalls every
possible assignment without error must satisfy
\begin{equation}
  Q\geq N\log_2\abs{\mathcal V},
  \qquad\text{equivalently}\qquad
  N\leq\frac{Q}{\log_2\abs{\mathcal V}}.
  \label{eq:finite-state-bound}
\end{equation}
\end{proposition}
\noindent\emph{Proof sketch.}
There are \(\abs{\mathcal V}^N\) possible independent value assignments.
Perfect recall requires that assignments needing different answers not collapse
to the same persistent state.  A \(Q\)-bit state has at most \(2^Q\) distinct
physical encodings, so injectivity requires
\(2^Q\geq\abs{\mathcal V}^N\).  Taking base-two logarithms yields
\eqref{eq:finite-state-bound}. \hfill\(\square\)

This counting bound is deliberately optimistic: it ignores floating-point
noise, hash collisions, readout interference, rank limits, replacement, and
optimization error.  The implemented consequences are concrete.  Exact-cache
records outside the stable top-\(C_x\) set are absent from that lane;
local-attention records older than \(W\) are absent from that lane; an
overwritten episodic capsule is absent entirely; and retained archive
numerators are generally lossy.  Increasing \(A\) extends
\eqref{eq:retained-horizon} without adding archive slots, but forces more
segments through the same rank-\(r\) capsule and can only increase the
opportunity for approximation and interference.  No equation in this section
supports exact recall over an unbounded stream.

\subsection{Code correspondence and validation boundary}
\label{subsec:code-correspondence}

The equations above were cross-checked against the following implementation
sites:
\begin{itemize}
  \item \texttt{gated\_delta.py}: decay, erase/write recurrence, RMS surprise,
        fp32 fast weights, and write-then-read order;
  \item \texttt{local\_attention.py}: append/truncate/read order, learned
        recency slope, fp32 softmax, and \(W\)-bounded state;
  \item \texttt{exact\_cache.py}: detached strict replacement, batch-local
        first-empty/minimum selection, exact softmax read, and mixed-dtype
        state accounting;
  \item \texttt{episodic.py}: positive feature map, separate \(M\) and \(z\),
        TensorSketch routing, stabilized mass normalization, ring lifecycle,
        retained horizon, and state formula;
  \item \texttt{compression.py}: reduced QR, small-core SVD, symmetric singular
        split, discarded tail, and additive error ledger;
  \item \texttt{model.py}: parallel branch execution, token-wise softmax
        fusion, residual ordering, and SwiGLU; and
  \item \texttt{state.py}: the exact persistent tensor fields counted in
        \eqref{eq:model-state-byte-bound}.
\end{itemize}

The corresponding tests check Gated Delta scan/step equality and fixed state;
local scan/step equality, window capping, and causality; exact-cache top-priority
retention, strict ties, scan/step equality, causality, dtype handling, and
constant allocation; episodic denominator separation, conditional routing
mass conservation, chunk equivalence through archive eviction, causality, and
bounded state; dense-versus-small-core compression and the error ledger; and
full-model chunk/step equality in evaluation mode.  The shipped release record
\texttt{results/release\_validation.json} reports fifty passing tests for the
audited snapshot.

Two measured artifacts are useful only as implementation anchors, not as
theorems.  The protocol
\texttt{aluclu\_diagnostic\_next\_token\_mqar\_v1} records \(90.43\%\)
answer accuracy and \(9{,}200\) bytes of batch-one state for its small learned
CPU diagnostic.  The protocol
\texttt{aluclu\_reference\_performance\_v2} records \(164{,}576\) bytes after
a saturated batch-one stream, a required warmup horizon of \(528\) tokens, and
a median \(143.71\) recurrent tokens/s in its reported single-thread CPU
environment.  These values establish that the specified mechanics executed
under those exact configurations; they do not prove large-scale language-model
quality, hardware-independent speed, or superiority over another architecture.
